\section{Geometry: the affine and isoparametric variants}
\label{sec:geom}

The reference element is the unit cube $[0,1]^3$ with trilinear shape functions. The Jacobian
$\Jmat(\xi) = \partial\xvec/\partial\xi$ is in general a function of position; the two variants
differ in whether they admit that.

\subsection{What each variant commits to}

The \emph{isoparametric} variant evaluates the geometry where it is needed, so $\Jmat$, its
adjugate and the surface areas all vary within the element
(\code{cvfem\_hex8\_geom\_at}, and the \code{\_isoparam} kernels).

The \emph{affine} variant evaluates $\Jmat$ once, at the reference centre
$\xi = (\tfrac12,\tfrac12,\tfrac12)$, and uses it everywhere
(\code{cvfem\_hex8\_affine\_adj}). For a parallelepiped this is exact. For a general trilinear
hexahedron it is an approximation, and it is the \emph{defining} approximation of the variant:
having made it, the variant's geometry \emph{is} that one constant Jacobian, and a kernel on this
path has no entitlement to any other geometric information. This is a stronger statement than it
first appears, and Section~\ref{sec:geom:edges} shows where ignoring it introduced a real bug.

What is stored per element is the adjugate and the determinant, not $\Jmat$ itself. With
$\adj(\mathsf{A}) = \det(\mathsf{A})\,\mathsf{A}^{-1}$ the code keeps
\code{cof0..cof8} $=\adj(\Jmat)$ row-major and \code{det} $=\det\Jmat$.

\subsection{Surface areas}

The physical area vector of surface $\scs$ in direction $q = \scs/4$ is the reference area vector
pushed forward by the cofactor matrix. Because $\Avec^{\mathrm{ref}}$ is $\tfrac14$ along axis
$q$, this reduces to a scaling of one row of the adjugate:
\begin{equation}
\Avec_\scs \;=\; \tfrac14 \bigl(\adj\Jmat\bigr)_{q,\cdot} ,
\end{equation}
which is exactly the nine lines \code{Ax0 = qtr*c0} \dots{} \code{Az2 = qtr*c8} in the geometry
loop of each sweep. All four surfaces of a direction group share one area vector, so the sweep
holds three of them, not twelve.

\subsection{Edge vectors, and why the Jacobian suffices}
\label{sec:geom:edges}

Several components need $\dvec_\scs$, the vector from node $i$ to node $j$ across the surface.
The obvious implementation differences the node coordinates, $\dvec_\scs = \xvec_j - \xvec_i$,
and that is what the code did until recently.

On the affine path it is the wrong choice, and not because it is inaccurate --- because it is
\emph{inconsistent}. Evaluating the trilinear Jacobian at the reference centre gives, for
direction $q$,
\begin{equation}
\bigl(\Jmat\bigr)_{\cdot,q} \;=\; \frac{1}{4}\!\!\sum_{\scs \,\in\, \text{group } q}\!\!
      \bigl(\xvec_{j(\scs)} - \xvec_{i(\scs)}\bigr),
\label{eq:col-is-mean}
\end{equation}
that is, column $q$ of the centre Jacobian is exactly the arithmetic mean of the four true edge
vectors of group $q$. This holds for \emph{any} trilinear hexahedron, not only a parallelepiped:
it is an identity of the trilinear basis evaluated at the centre, and it was verified
symbolically against the code's own adjugate expressions on cubes, sheared parallelepipeds and
randomly distorted elements, agreeing to $10^{-16}$ in every case.

So the affine variant has already committed to that mean when it computes the flux. A kernel that
then differences node coordinates for $\dvec$ is using a \emph{second, different} geometry for one
term while the rest of the same kernel uses the first. On a cube the two coincide exactly; on a
distorted element the true edges differ from their mean by $10\%$ to over $100\%$, which is a
measure of how far outside its validity the affine variant has been pushed, not of an error in
either formula.

The affine path therefore takes
\begin{equation}
\dvec_\scs \;=\; \bigl(\Jmat\bigr)_{\cdot,\,\scs/4},
\end{equation}
and recovers $\Jmat$ from what is already in registers using the $3\times 3$ identity
\begin{equation}
\adj\bigl(\adj \mathsf{A}\bigr) \;=\; \det(\mathsf{A})\,\mathsf{A}
\qquad\Longrightarrow\qquad
\Jmat \;=\; \frac{\adj(\adj\Jmat)}{\det\Jmat} ,
\label{eq:adj-adj}
\end{equation}
implemented in \code{cvfem\_hex8\_affine\_edge\_cols}. The cofactors and $1/\det$ are already live
in the geometry loop that builds the area vectors, so the three columns cost roughly forty flops
per element and no loads, and the four surfaces of a group share one column. The isoparametric
kernels keep true per-surface coordinates, which is correct where the geometry genuinely varies.

\paragraph{A consistency obligation this creates.} Every \emph{producer} of precomputed data must
use the same $\dvec$ as every \emph{consumer}. The coefficient table
(\code{cvfem\_hex8\_build\_rc\_coeff}) and the partially assembled tangent
(\code{cvfem\_hex8\_build\_pa\_tangent}) are both producers, and both had to move to
Eq.~\eqref{eq:adj-adj} with the kernels. Missing the second of them put the stored tangent and the
kernel reading it on different geometries, which \code{cvfem\_pa\_tangent\_test} caught at
$1.4\times10^{-5}$ --- five orders of magnitude above round-off, and invisible to every
cube-based check, because on a cube Eq.~\eqref{eq:col-is-mean} makes the two definitions
identical. Any future change here needs a non-affine oracle.
